
\chapter{עבודות קשורות}

\section{חילוץ אותות בשיטות קלאסיות}

גישות קלאסיות לחילוץ אותות סינוסואידליים כוללות: \textenglish{MUSIC} (\textenglish{MUltiple SIgnal Classification}), \textenglish{ESPRIT} (\textenglish{Estimation of Signal Parameters via Rotational Invariance Techniques}) ו-\textenglish{CLEAN}. שיטות אלו מבוססות על פירוק ערכים עצמיים (\textenglish{eigendecomposition}) של מטריצת הקורלציה ומציגות ביצועים מצוינים בתנאי \textenglish{SNR} גבוה, אך רגישות לסטיות ממודל האות הנחות \cite{williamson2016complex}.

\section{גישות למידה עמוקה לעיבוד אותות}

\textenglish{WaveNet} \cite{engel2017wavenet} הציג ארכיטקטורת קונבולוציה מורחבת (\textenglish{dilated convolution}) לסינתזת אות. \textenglish{Conv-TasNet} \cite{luo2019convtasnet} הרחיב את הגישה להפרדת דוברים, תוך שהוא מחליף את ה-\textenglish{STFT} בקנה מידה למידה (\textenglish{learned encoder-decoder}). \textenglish{sudo rm-rf} \cite{tzinis2020universal} הציג מנגנון \textenglish{U-Net} סדרתי לפריד אותות כלליים.

\section{שיטות משולבות עם מנגנוני קשב}

שילוב \textenglish{attention} עם \textenglish{RNN} לעיבוד אותות החל ב-\cite{nachmani2020voice}, שם הוכח שמנגנון קשב "ממוקד דוברים" (\textenglish{speaker-aware attention}) משפר הפרדת קולות. גישות \textenglish{Transformer} לעיבוד אותות \cite{vaswani2017attention} דורשות עיבוד מקביל מלא ולכן פחות מתאימות לאינפרנס בזמן אמת, בניגוד לגישת ה-\textenglish{BiLSTM} שלנו.

\section{מדדי הערכה בתחום}

מדד ה-\textenglish{SI-SNR} אומץ כסטנדרט הזהב בתחום הפרדת אותות לאחר מאמר ה-\textenglish{tasnet} \cite{luo2019convtasnet}. מדדים נוספים כמו \textenglish{PESQ} (\textenglish{Perceptual Evaluation of Speech Quality}) ו-\textenglish{STOI} (\textenglish{Short-Time Objective Intelligibility}) משמשים בהקשרי שיפור דיבור \cite{valentini2019speech}.

\section{השוואה מקיפה}

בהשוואה לכל גישות הבסיס שניסינו, ה-\textenglish{BiLSTM} שלנו מציג שיפור עקבי ומובהק סטטיסטית (\textenglish{p < 0.01}, מבחן \textenglish{Wilcoxon signed-rank}) על כל מערכי הנתונים שנבדקו. יחד עם זאת, מודגש כי \textenglish{Conv-TasNet} \cite{luo2019convtasnet} הנו מתחרה קרוב ויש לו יתרון מבחינת \textenglish{latency} נמוכה יותר (\textenglish{$\sim$12 ms} לעומת \textenglish{$\sim$35 ms} שלנו).
